\chapter{The Viscous Hubble Diagram and Distance Modulus}
\label{app:hubble_diagram}
\addtocontents{toc}{\protect\setcounter{tocdepth}{0}}

\noindent\textbf{Abstract:} Standard cosmology evaluates the luminosity distance of distant supernovae through the expansion history of the FLRW metric. Tension Field Theory instead derives a nonlinear coordinate-distance law from coherent-action-preserving visco-elastic frequency relaxation. Under an additional macroscopic source-envelope timing closure, this yields a conditional TFT luminosity-distance relation and a directly testable non-expanding Hubble diagram. This appendix also states the Euclidean acoustic-projection closure used in Part III.

\section{The Two Factors of Viscous Luminosity Distance ($d_{\mathrm{L}}$)}

As derived in \hyperref[app:photon_dispersion_relaxation]{Appendix~\ref*{app:photon_dispersion_relaxation}}, the general TFT redshift law is
\begin{equation}
    \ln(1+z)
    =
    \int_{\mathrm{path}}
    \frac{\Gamma_{\mathrm{D}}(\ell)}
    {v_{\mathrm{ray}}\!\left[\omega(\ell),\alpha_{\mathrm{G}}(\ell)\right]}
    \,\mathrm{d}\ell .
\end{equation}

In the weak-loading, constant-rate Euclidean limit,
\begin{equation}
    \alpha_{\mathrm{G}}\to0,
    \qquad
    v_{\mathrm{ray}}\simeq c_{\mathrm{E}},
    \qquad
    \Gamma_{\mathrm{D}}\simeq\Gamma_{\mathrm{D}}^{(\mathrm{cal})},
\end{equation}
the coordinate distance becomes
\begin{equation}
    d_{\mathrm{coord}}
    =
    \frac{c_{\mathrm{E}}}{\Gamma_{\mathrm{D}}^{(\mathrm{cal})}}
    \ln(1+z).
\end{equation}

Using the empirical low-redshift calibration
\begin{equation}
    \Gamma_{\mathrm{D}}^{(\mathrm{cal})}
    =
    H_0^{(\mathrm{obs})},
\end{equation}
this may be written
\begin{equation}
    d_{\mathrm{coord}}
    =
    \frac{c_{\mathrm{E}}}{H_0^{(\mathrm{obs})}}
    \ln(1+z).
\end{equation}

The observational luminosity distance requires a second transport statement. The two relevant factors are:

\begin{itemize}
    \item \textbf{Factor A (Coherent Energy Relaxation):} Coherent-action preservation gives
    \begin{equation}
        \frac{E_{\mathrm{obs}}}{E_{\mathrm{emit}}}
        =
        \frac{\omega_{\mathrm{obs}}}{\omega_{\mathrm{emit}}}
        =
        \frac{1}{1+z}.
    \end{equation}

    \item \textbf{Factor B (Conditional Macroscopic Arrival-Rate Dilation):} For a fixed-cycle coherent packet,
    \begin{equation}
        \Delta t_{\mathrm{pkt,obs}}
        =
        (1+z)\Delta t_{\mathrm{pkt,emit}}.
    \end{equation}
    As emphasized in \hyperref[chap:cosmic_redshift]{Chapter~\ref*{chap:cosmic_redshift}}, this packet result does not automatically establish the dilation of an arbitrary astronomical source envelope. If the macroscopic source-to-envelope transport closure inherits the same $(1+z)$ arrival-rate factor, then the received bolometric power satisfies
    \begin{equation}
        \frac{P_{\mathrm{obs}}}{P_{\mathrm{emit}}}
        =
        \frac{1}{(1+z)^2}.
    \end{equation}
\end{itemize}

Under this macroscopic envelope-dilation closure, the inverse-square flux relation converts the two power factors into one multiplicative factor of $(1+z)$ in the effective luminosity distance:
\begin{equation}
    d_{\mathrm{L(TFT)}}
    =
    d_{\mathrm{coord}}(1+z).
\end{equation}

Therefore,
\begin{equation}
    d_{\mathrm{L(TFT)}}
    =
    \frac{c_{\mathrm{E}}}
    {\Gamma_{\mathrm{D}}^{(\mathrm{cal})}}
    (1+z)\ln(1+z),
\end{equation}
or, using the present empirical calibration,
\begin{equation}
    d_{\mathrm{L(TFT)}}
    =
    \frac{c_{\mathrm{E}}}
    {H_0^{(\mathrm{obs})}}
    (1+z)\ln(1+z).
\end{equation}

This luminosity-distance law is therefore a conditional TFT prediction: the coordinate-distance component follows from the coherent-relaxation sector, whereas the additional $(1+z)$ factor requires completion of the macroscopic source-envelope transport law.

\section{The Distance Modulus}

Astronomical observations commonly express luminosity distance through the distance modulus,
\begin{equation}
    \mu = m-M
    =
    5\log_{10}
    \left(
        \frac{d_{\mathrm{L}}}{\mathrm{Mpc}}
    \right)
    +25.
\end{equation}

Substituting the conditional TFT luminosity distance gives
\begin{equation}
    \mu_{\mathrm{TFT}}
    =
    5\log_{10}
    \left[
        \frac{c_{\mathrm{E}}}
        {H_0^{(\mathrm{obs})}\,\mathrm{Mpc}}
        (1+z)\ln(1+z)
    \right]
    +25.
\end{equation}

This expression defines the TFT Hubble-diagram prediction to be tested directly against supernova distance-modulus data under the stated constant-rate and macroscopic envelope-dilation closures.

\section{Taylor Expansion and Cosmographic Comparison}

The exact nonlinear factor appearing in the conditional TFT luminosity distance has the Taylor expansion
\begin{equation}
    (1+z)\ln(1+z)
    =
    z
    +
    \frac{z^2}{2}
    -
    \frac{z^3}{6}
    +
    \frac{z^4}{12}
    +\mathcal{O}(z^5).
\end{equation}

Therefore,
\begin{equation}
    d_{\mathrm{L(TFT)}}
    =
    \frac{c_{\mathrm{E}}}{H_0^{(\mathrm{obs})}}
    \left[
        z
        +
        \frac{z^2}{2}
        -
        \frac{z^3}{6}
        +
        \mathcal{O}(z^4)
    \right].
\end{equation}

At low redshift, the leading term recovers the linear Hubble relation. For comparison, the standard cosmographic luminosity-distance expansion begins as
\begin{equation}
    d_{\mathrm{L}}
    \simeq
    \frac{c_{\mathrm{E}}}{H_0}
    \left[
        z
        +
        \frac{1-q_0}{2}z^2
        +\cdots
    \right].
\end{equation}

Matching only the quadratic coefficient of the TFT expansion gives
\begin{equation}
    \frac{1-q_0}{2}
    =
    \frac{1}{2}
    \qquad\Longrightarrow\qquad
    q_0=0.
\end{equation}

Consequently, the positive $\frac{1}{2}z^2$ term is not by itself a mathematical derivation of the negative-$q_0$ acceleration inferred within standard cosmography. TFT instead proposes a distinct non-expanding nonlinear distance law. Its ability to reproduce the observed supernova Hubble diagram must therefore be assessed by fitting the complete expression
\begin{equation}
    d_{\mathrm{L(TFT)}}
    =
    \frac{c_{\mathrm{E}}}{H_0^{(\mathrm{obs})}}
    (1+z)\ln(1+z)
\end{equation}
directly to the observational data rather than identifying individual Taylor coefficients with FLRW kinematic parameters.

\section{Euclidean Acoustic Projection Closure}
\label{sec:euclidean_acoustic_projection_closure}

The characteristic Viscoelastic Vacuum Relaxation length is
\begin{equation}
    R_{\mathrm{U}}
    =
    \frac{c_{\mathrm{E}}}
    {\Gamma_{\mathrm{D}}^{(\mathrm{cal})}}.
\end{equation}

As established in \hyperref[app:photon_dispersion_relaxation]{Appendix~\ref*{app:photon_dispersion_relaxation}}, $R_{\mathrm{U}}$ is the characteristic one-$e$-fold coherent-relaxation length. It is not a universal photon stopping radius, and the frequency-dependent CMB assimilation threshold may occur at substantially larger distances.

For the Euclidean reinterpretation of the observed CMB acoustic angular scale developed in \hyperref[chap:cosmological_anomalies]{Chapter~\ref*{chap:cosmological_anomalies}}, TFT introduces the additional constitutive projection closure
\begin{equation}
    D_{\mathrm{proj}}
    \equiv
    R_{\mathrm{U}}.
\end{equation}

This identification is presently a cosmological geometric closure rather than a microscopic consequence of the photon-relaxation derivation. Under this closure, the effective TFT acoustic scale associated with an observed angular scale $\theta_*$ is
\begin{equation}
    r_{\mathrm{s(TFT)}}
    =
    D_{\mathrm{proj}}\theta_*
    =
    R_{\mathrm{U}}\theta_*.
\end{equation}

Using
\begin{equation}
    R_{\mathrm{U}}
    \approx
    4104.498\ \mathrm{Mpc},
    \qquad
    \theta_*
    \approx
    0.01041,
\end{equation}
gives
\begin{equation}
    r_{\mathrm{s(TFT)}}
    \approx
    42.73\ \mathrm{Mpc}.
\end{equation}

The dimensional consistency is immediate:
\begin{equation}
    [D_{\mathrm{proj}}]
    =
    [R_{\mathrm{U}}]
    =
    [r_{\mathrm{s(TFT)}}]
    =
    \mathrm{m},
    \qquad
    [\theta_*]=1.
\end{equation}

The numerical $42.73$ Mpc scale is therefore a well-defined consequence of the proposed closure $D_{\mathrm{proj}}=R_{\mathrm{U}}$, but it is not yet an independent parameter-free prediction. Closing this branch requires a microscopic or percolative derivation showing why the relevant CMB projection distance must coincide with the one-$e$-fold vacuum-relaxation scale.
